\section*{Розділ II. Організація та скликання КСУ}
\addcontentsline{toc}{section}{Розділ II. Організація та скликання КСУ}

\subsection*{2.1. Порядок скликання чергової КСУ}
\addcontentsline{toc}{subsection}{2.1. Порядок скликання чергової КСУ}
    2.1.1. Чергова КСУ скликається Спікером КСУ не рідше одного разу на рік відповідно до пункту 3.1.2 Положення про органи студентського самоврядування Київського авіаційного інституту.

    2.1.2. Спікер КСУ зобов'язаний скликати чергову КСУ, якщо з дня попереднього засідання минуло більше 11 місяців.

    2.1.3. Чергова КСУ скликається для розгляду планових питань діяльності органів студентського самоврядування, включаючи заслуховування звітів, затвердження кошторису та інших питань, визначених пунктом 3.2 Положення про ОСС.

\subsection*{2.2. Порядок скликання позачергової КСУ}
\addcontentsline{toc}{subsection}{2.2. Порядок скликання позачергової КСУ}
    2.2.1. Позачергова КСУ може бути скликана за ініціативою суб'єктів, визначених пунктом 3.1.2 Положення про ОСС: Студентської ради Київського авіаційного інституту (СР КАІ), Студентської ради студмістечка (СР СМ), Студентської уповноваженої делегації (СУД), Центральної виборчої комісії студентів (ЦВКс), Спікера КСУ, не менше третини від загального складу делегатів КСУ або за ініціативою не менше 5\% студентів від загальної кількості студентів Університету.

    2.2.2. Ініціатива про скликання позачергової КСУ подається Спікеру КСУ у письмовій формі з обов'язковим зазначенням питань, що пропонуються до розгляду, та обґрунтуванням необхідності позачергового скликання.

    2.2.3. Спікер КСУ зобов'язаний скликати позачергову КСУ протягом 15 календарних днів з дня отримання відповідної ініціативи, якщо вона відповідає вимогам цього Регламенту та Положення про ОСС.

\subsection*{2.3. Обов'язки та повноваження Спікера КСУ}
\addcontentsline{toc}{subsection}{2.3. Обов'язки та повноваження Спікера КСУ}
    2.3.1. Спікер КСУ виконує організаційні функції щодо роботи КСУ відповідно до пункту 3.4.4 Положення про ОСС та цього Регламенту.

    2.3.2. До основних обов'язків Спікера КСУ належить:

        \begin{enumerate}[label=\alph*)]
            \item Скликання засідань КСУ у встановленому порядку;
            \item Ведення засідань КСУ;
            \item Підготовка проєкту порядку денного засідання КСУ на основі отриманих пропозицій;
            \item Забезпечення дотримання цього Регламенту під час засідань КСУ;
            \item Координація взаємодії з органами студентського самоврядування з питань підготовки засідань КСУ;
            \item Підписання протоколів засідань КСУ;
            \item Забезпечення оприлюднення рішень КСУ у встановленому порядку.
        \end{enumerate}

    2.3.3. Спікер КСУ має право:

        \begin{enumerate}[label=\alph*)]
            \item Визначати технічні аспекти проведення засідання (платформа для дистанційної участі, порядок реєстрації учасників тощо);
            \item Запрошувати на засідання КСУ представників адміністрації Університету, експертів та інших осіб для надання інформації з питань порядку денного;
            \item Призупиняти засідання у разі порушення порядку або технічних проблем.
        \end{enumerate}

\subsection*{2.4. Повідомлення про скликання КСУ}
\addcontentsline{toc}{subsection}{2.4. Повідомлення про скликання КСУ}
    2.4.1. Повідомлення про скликання чергової КСУ надсилається делегатам КСУ не пізніше ніж за 7 (сім) календарних днів до дати засідання.

    2.4.2. Повідомлення про скликання позачергової КСУ надсилається делегатам КСУ не пізніше ніж за 3 (три) календарні дні до дати засідання.

    2.4.3. Повідомлення про скликання КСУ повинно містити:

        \begin{enumerate}[label=\alph*)]
            \item Дату, час та спосіб проведення засідання (очно, дистанційно або в змішаному форматі);
            \item Проєкт порядку денного;
            \item Матеріали до питань порядку денного (за наявності);
            \item Технічну інформацію для участі у засіданні (посилання на платформу, інструкції для підключення тощо).
        \end{enumerate}

    2.4.4. Повідомлення надсилається офіційними каналами зв'язку органів студентського самоврядування та розміщується на їх офіційних інформаційних ресурсах.

\subsection*{2.5. Підготовка засідання КСУ}
\addcontentsline{toc}{subsection}{2.5. Підготовка засідання КСУ}
    2.5.1. Проєкт порядку денного засідання КСУ формується Спікером КСУ на основі:

        \begin{enumerate}[label=\alph*)]
            \item Пропозицій від органів студентського самоврядування;
            \item Питань, обов'язкових для розгляду відповідно до Положення про ОСС;
            \item Ініціатив делегатів КСУ;
            \item Питань, що випливають з рішень попередніх засідань КСУ.
        \end{enumerate}

    2.5.2. Пропозиції до порядку денного подаються Спікеру КСУ у письмовій формі не пізніше ніж за 10 календарних днів до планованої дати засідання для чергової КСУ та не пізніше ніж за 5 календарних днів для позачергової КСУ.

    2.5.3. Спікер КСУ забезпечує підготовку матеріалів до засідання, включаючи:

        \begin{enumerate}[label=\alph*)]
            \item Проєкти рішень з питань порядку денного;
            \item Супровідні документи та довідкові матеріали;
            \item Технічне забезпечення проведення засідання.
        \end{enumerate}

\subsection*{2.6. Форми проведення засідань КСУ}
\addcontentsline{toc}{subsection}{2.6. Форми проведення засідань КСУ}
    2.6.1. Засідання КСУ може проводитися у таких формах:

        \begin{enumerate}[label=\alph*)]
            \item Дистанційно -- з використанням засобів електронного зв'язку (відеоконференції, системи електронного голосування тощо);
            \item Очно -- з фізичною присутністю делегатів у визначеному місці;
            \item У змішаному форматі -- з можливістю участі як очно, так і дистанційно.
        \end{enumerate}

    2.6.2. Форма проведення засідання визначається Спікером КСУ з урахуванням практичних можливостей та обставин.

    2.6.3. При проведенні засідання КСУ дистанційно або у змішаному форматі забезпечується:

        \begin{enumerate}[label=\alph*)]
            \item Належна ідентифікація учасників засідання;
            \item Можливість вільного обговорення питань порядку денного;
            \item Достовірна фіксація результатів голосування;
            \item Ведення протоколу засідання.
        \end{enumerate}

    2.6.4. Технічні вимоги та процедури проведення дистанційних засідань можуть деталізуватися окремими рішеннями КСУ або методичними рекомендаціями Спікера КСУ.